%=======================================================================
% mmae450/tex/lectures/lecture_fvm_1D_intro.tex
% Standalone Beamer deck that inputs an existing Chapter 5 .tikz figure
% located in: mmae450/tex/figures/lectures/
%
% IMPORTANT:
% The .tikz file itself does: \input{tex/preamble/tikz_house_styles.tex}
% so we add the repo root to TeX's input search path using \input@path.
%=======================================================================


\documentclass[aspectratio=169]{beamer}

% --- Your beamer styling (adjust path/name if yours differs)
\input{../preamble/beamer_styles.tex}

% --- TikZ (beamer_styles.tex may already load it; harmless if duplicated)
\usepackage{tikz}
\usetikzlibrary{arrows.meta,calc,positioning,decorations.pathreplacing}

% --- Make sure TeX can resolve paths like tex/preamble/... used INSIDE the .tikz
% From mmae450/tex/lectures, the repo root is ../../
\makeatletter
\def\input@path{{../../}{../}{./}}
\makeatother

\title{Finite Volume Method in 1D}
\subtitle{Finite Volume Method for Incompressible Fluid Flow}
\author{MMAE 450}
\date{}

\begin{document}

% ============================================================
% CONSISTENCY UPDATE (Projection Method Clarification)
% ============================================================
% Time discretization used throughout this deck:
%
%   (v^{n+1} - v^n)/dt = N(v^n) - (1/rho) grad p^{n+1}
%
% where:
%   N(v^n) = - (v · grad)v + nu laplacian(v)   evaluated at time n.
%
% Predictor:
%   v^* = v^n + dt * N(v^n)
%
% Corrector:
%   v^{n+1} = v^* - (dt/rho) grad p^{n+1}
%
% All convective and diffusive fluxes are evaluated at time level n (explicit).
% Pressure gradient is evaluated at time level n+1 (implicit).
%
% Finite Volume Poisson equation:
%
%   sum_f ( (dp/dn)_f l_f ) = (rho/dt) sum_f phi_f^*
%
% which on a Cartesian grid becomes:
%
%   laplacian(p) = (rho/dt) div(v^*)
%
% This matches the implementation in the companion notebook.
% ============================================================


% ----------------------------------------------------------------------
\begin{frame}
  \titlepage
\end{frame}

%==================== Slide 1 ====================
\begin{frame}{Incompressible Navier--Stokes (strong form)}
\textbf{Unknowns:} velocity $\mathbf{v}(\mathbf{x},t)$ and pressure $p(\mathbf{x},t)$.

\bigskip
\textbf{Continuity (incompressibility):}
\[
\nabla\cdot\mathbf{v}=0
\]

\textbf{Momentum:}
\[
\rho\left(\frac{\partial \mathbf{v}}{\partial t}+(\mathbf{v}\cdot\nabla)\mathbf{v}\right)
=
-\nabla p+\mu\nabla^2\mathbf{v}
\]
\end{frame}
%==================== Slide 2 ====================
\begin{frame}{Rewrite convection in conservative (flux) form}
For incompressible flow, $\nabla\cdot\mathbf{v}=0$, so
\[
(\mathbf{v}\cdot\nabla)\mathbf{v}
=
\nabla\cdot(\mathbf{v}\otimes\mathbf{v})
\]
(where $(\mathbf{v}\otimes\mathbf{v})_{ij}=v_i v_j$).

\bigskip
Then the momentum equation becomes
\[
\rho\frac{\partial \mathbf{v}}{\partial t}
+
\nabla\cdot\big(\rho\,\mathbf{v}\otimes\mathbf{v}\big)
=
-\nabla p+\mu\nabla^2\mathbf{v}.
\]
\end{frame}

%==================== Slide 3 ====================
\begin{frame}{Put everything into $\nabla\cdot(\cdot)$ form}
Use the identities
\[
-\nabla p = -\nabla\cdot(p\mathbf{I}),
\qquad
\mu\nabla^2\mathbf{v} = \nabla\cdot(\mu\nabla\mathbf{v})
\quad (\mu=\text{const}).
\]

\bigskip
So we can write
\[
\rho\frac{\partial \mathbf{v}}{\partial t}
+
\nabla\cdot\Big(\rho\,\mathbf{v}\otimes\mathbf{v}
+ p\mathbf{I}
-\mu\nabla\mathbf{v}\Big)
=\mathbf{0}.
\]
\end{frame}
%==================== Slide 4 ====================
\begin{frame}{Flux form for finite volume}
Define the \textbf{momentum flux tensor}
\[
\mathbf{F}
\;\equiv\;
\rho\,\mathbf{v}\otimes\mathbf{v}
+ p\mathbf{I}
-\mu\nabla\mathbf{v}.
\]

\bigskip
Then the momentum equation is
\[
\boxed{
\rho \frac{\partial \mathbf{v}}{\partial t}
+\nabla\cdot\mathbf{F}
=
\mathbf{0}.
}
\]

\bigskip
\textbf{Interpretation:} \emph{storage} $+$ \emph{net outward flux} $=0$ (per unit volume).
\end{frame}

%==================== Slide 5 ====================
\begin{frame}{Finite volume idea: integrate over a 2D control area}
Consider a single 2D finite volume cell (per unit thickness) with area $A_P$
and boundary $\partial A_P$.

\bigskip
Start from the flux form
\[
\rho \frac{\partial \mathbf{v}}{\partial t}+\nabla\cdot\mathbf{F}=\mathbf{0}.
\]

\bigskip
Integrate over the control area $A_P$:
\[
\int_{A_P}\rho \frac{\partial \mathbf{v}}{\partial t}\,dA
+
\int_{A_P}\nabla\cdot\mathbf{F}\,dA
=
\mathbf{0}.
\]

\bigskip
\textbf{Goal:} convert the divergence term to a \emph{boundary flux} integral.
\end{frame}
%==================== Slide 6 ====================
\begin{frame}{Divergence theorem in 2D (per unit thickness)}
Apply the divergence theorem:
\[
\int_{A_P}\nabla\cdot\mathbf{F}\,dA
=
\int_{\partial A_P}\mathbf{F}\cdot\mathbf{n}\,ds,
\]
where $\mathbf{n}$ is the outward unit normal on $\partial A_P$ and $ds$ is
arc length along the boundary.

\bigskip
So the finite-volume integral statement is
\[
\boxed{
\int_{A_P}\rho \frac{\partial \mathbf{v}}{\partial t}\,dA
+
\int_{\partial A_P}\mathbf{F}\cdot\mathbf{n}\,ds
=
\mathbf{0}.
}
\]

\bigskip
\textbf{Interpretation:} \emph{rate of change of momentum in the cell}
$+$ \emph{net outward momentum flux through the cell boundary} $=0$.
\end{frame}

%==================== Slide 7 ====================
\begin{frame}{From boundary integral to a sum over faces}
Approximate the boundary integral by summing fluxes over the cell faces $f$:
\[
\int_{\partial A_P}\mathbf{F}\cdot\mathbf{n}\,ds
\;\approx\;
\sum_{f\in\partial P}\left(\mathbf{F}_f\cdot\mathbf{n}_f\right)\,\ell_f,
\]
where $\ell_f$ is the length of face $f$ (per unit thickness) and
$\mathbf{n}_f$ is the outward unit normal on that face.

\bigskip
Also approximate the storage term using a cell-centered value $\mathbf{v}_P(t)$:
\[
\int_{A_P}\rho \frac{\partial \mathbf{v}}{\partial t}\,dA
\;\approx\;
\rho\,A_P\,\frac{d\mathbf{v}_P}{dt}.
\]

\bigskip
\textbf{Finite-volume semi-discrete form:}
\[
\boxed{
\rho\,A_P\,\frac{d\mathbf{v}_P}{dt}
+
\sum_{f\in\partial P}\left(\mathbf{F}_f\cdot\mathbf{n}_f\right)\,\ell_f
=
\mathbf{0}.
}
\]
\end{frame}

%==================== Slide 8 ====================
\begin{frame}{Decomposing the face flux $\mathbf{F}_f\cdot\mathbf{n}_f$}
Recall the definition of the momentum flux tensor:
\[
\mathbf{F}
=
\rho\,\mathbf{v}\otimes\mathbf{v}
+
p\mathbf{I}
-
\mu\nabla\mathbf{v}.
\]

\bigskip
At a face $f$, the normal flux becomes
\[
\mathbf{F}_f\cdot\mathbf{n}_f
=
\underbrace{\rho\,\mathbf{v}_f(\mathbf{v}_f\cdot\mathbf{n}_f)}_{\text{convection}}
+
\underbrace{p_f\,\mathbf{n}_f}_{\text{pressure}}
-
\underbrace{\mu\,(\nabla\mathbf{v})_f\cdot\mathbf{n}_f}_{\text{diffusion}}.
\]

\bigskip
\textbf{So each face contributes:}
\[
(\text{convective flux})
+
(\text{pressure force})
+
(\text{viscous traction}).
\]

\bigskip
All physics enters \emph{through faces}.
\end{frame}

%==================== Slide 9 ====================
\begin{frame}{Introducing the normal flux $\phi_f$}
Define the \textbf{normal volumetric flux} through face $f$:
\[
\phi_f
\;\equiv\;
(\mathbf{v}_f \cdot \mathbf{n}_f)\,\ell_f.
\]

\bigskip
Then the convective momentum flux becomes
\[
\rho\,\mathbf{v}_f(\mathbf{v}_f\cdot\mathbf{n}_f)\,\ell_f
=
\rho\,\mathbf{v}_f\,\phi_f.
\]

\bigskip
So the semi-discrete FV equation can be written as
\[
\rho A_P \frac{d\mathbf{v}_P}{dt}
+
\sum_{f\in\partial P}
\Big(
\rho\,\mathbf{v}_f\,\phi_f
+
p_f\,\mathbf{n}_f\,\ell_f
-
\mu\,(\nabla\mathbf{v})_f\cdot\mathbf{n}_f\,\ell_f
\Big)
=
\mathbf{0}.
\]

\bigskip
\textbf{Key idea:} incompressibility will require
\[
\sum_{f\in\partial P} \phi_f = 0.
\]
\end{frame}

%==================== Slide 10 ====================
\begin{frame}{Key idea: incompressibility in finite-volume form}
For an incompressible fluid,
\[
\nabla \cdot \mathbf{v} = 0.
\]

\bigskip
Integrate over a control area $A_P$:
\[
\int_{A_P} \nabla \cdot \mathbf{v}\, dA = 0.
\]

\bigskip
Apply the divergence theorem:
\[
\int_{\partial A_P} \mathbf{v}\cdot \mathbf{n}\, ds = 0.
\]

\bigskip
Discrete (finite-volume) form:
\[
\boxed{
\sum_{f\in\partial P} \phi_f = 0.
}
\]

\end{frame}

\begin{frame}
\bigskip
\textbf{Interpretation:}

\begin{itemize}
\item No net volumetric flux leaves or enters the cell.
\item Fluid cannot accumulate or disappear inside $A_P$.
\item Pressure will be computed to enforce this constraint.
\end{itemize}
\end{frame}

%==================== Slide 11 ====================
\begin{frame}{Time discretization: predictor step}
Starting from the semi-discrete FV form
\[
\rho A_P \frac{d\mathbf{v}_P}{dt}
+
\sum_{f\in\partial P}
\Big(
\rho\,\mathbf{v}_f\,\phi_f
+
p_f\,\mathbf{n}_f\,\ell_f
-
\mu\,(\nabla\mathbf{v})_f\cdot\mathbf{n}_f\,\ell_f
\Big)
=
\mathbf{0},
\]

\bigskip
we advance in time using a simple forward Euler scheme.

\bigskip
\textbf{Predictor: ignore pressure for the moment}
\[
\rho A_P
\frac{\mathbf{v}_P^{*}-\mathbf{v}_P^{n}}{\Delta t}
+
\sum_{f\in\partial P}
\Big(
\rho\,\mathbf{v}_f^{n}\,\phi_f^{n}
-
\mu\,(\nabla\mathbf{v})_f^{n}\cdot\mathbf{n}_f\,\ell_f
\Big)
=
\mathbf{0}.
\]

\bigskip
This produces a provisional velocity $\mathbf{v}^{*}$ that generally
\emph{does not satisfy}
\[
\sum_{f\in\partial P} \phi_f^{*} = 0.
\]
\end{frame}

%==================== Slide 12 ====================
\begin{frame}{Implementation details: computing face quantities}
To evaluate the predictor step, we must approximate face values.

\bigskip
\textbf{1. Face velocity $\mathbf{v}_f$} ($\mathbf{v}_N$ is the velocity vector a the center of the neighboring cell).

\[
\mathbf{v}_f \approx \dfrac{\mathbf{v}_P + \mathbf{v}_N}{2} & \text{(central)}
\]
\textbf{2. Normal volumetric flux}

\[
\phi_f
=
(\mathbf{v}_f \cdot \mathbf{n}_f)\,\ell_f
\]

\bigskip
\textbf{3. Diffusive term}

\[
(\nabla \mathbf{v})_f \cdot \mathbf{n}_f
\approx
\frac{\mathbf{v}_N - \mathbf{v}_P}{d_{PN}}
\]

where $d_{PN}$ is the distance between neighboring cell centers.

\end{frame}
%==================== Slide 13 =======================
\begin{frame}{Diffusion term: normal gradient at a face}
\centering
\renewcommand{\FigWidth}{0.85\textwidth}
%=======================================================================
% chap05_figXX_2D_fv_diffusion_face.tikz
% 2D finite volume diffusion term at a face (HOUSE STYLE)
%=======================================================================

\providecommand{\FigWidth}{8.0cm}   % fallback only
\def\FigDesignWidth{8.0}           % includes label margins

\input{../preamble/tikz_house_styles.tex}

\begin{tikzpicture}[house]

  % ---- Styles ----
  \tikzset{
    halfMain/.style = {font=\figMainFont},
    halfAxis/.style = {font=\figAxisFont},
    axis/.style     = {axisLine, line width=0.9pt},
    face/.style     = {axisLine, line width=1.0pt},
    cellEdge/.style = {axisLine, line width=0.8pt},
    centerDot/.style= {fill=black},
    normalArrow/.style = {axisArrow, line width=0.9pt},
    brace/.style    = {decorate, decoration={brace, amplitude=4pt}, line width=0.6pt}
  }

  % ---- Layout parameters (DESIGN UNITS) ----
  \def\cellW{3.2}
  \def\cellH{2.2}
  \def\xP{0.5*\cellW}
  \def\xN{1.5*\cellW}
  \def\yC{0.5*\cellH}

  % ---- Cells P and N ----
  \draw[cellEdge] (0,0) rectangle (\cellW,\cellH);
  \draw[cellEdge] (\cellW,0) rectangle (2*\cellW,\cellH);

  % ---- Shared face ----
  \draw[face] (\cellW,0) -- (\cellW,\cellH);

  % ---- Cell centers ----
  \fill[centerDot] (\xP,\yC) circle (2pt);
  \fill[centerDot] (\xN,\yC) circle (2pt);

  \node[halfMain, below=3pt] at (\xP,\yC) {$P$};
  \node[halfMain, below=3pt] at (\xN,\yC) {$N$};

  % ---- Velocity labels ----
  \node[halfMain] at (\xP,0.78*\cellH) {$\mathbf{v}_P$};
  \node[halfMain] at (\xN,0.78*\cellH) {$\mathbf{v}_N$};

  % ---- Face midpoint ----
  \fill[centerDot] (\cellW,\yC) circle (1.6pt);
  \node[halfMain, right=3pt] at (\cellW,\yC+0.2) {$f$};

  % ---- Normal vector ----
  \draw[normalArrow] (\cellW,\yC) -- ++(1.0,0)
      node[halfMain, right] {$\mathbf{n}_f$};

  % ---- Distance between centers ----
  \draw[brace] (\xP,0.12*\cellH) -- (\xN,0.12*\cellH)
      node[midway, below=8pt, halfAxis] {$d_{PN}$};

  % ---- Diffusion approximation equation ----
  \node[halfMain, align=center] at (\cellW,-1.2)
  {$
  (\nabla \mathbf{v})_f \cdot \mathbf{n}_f
  \;\approx\;
  \dfrac{\mathbf{v}_N - \mathbf{v}_P}{d_{PN}}
  $};

\end{tikzpicture}
\end{frame}
%======================================================
\begin{frame}
\bigskip
\textbf{Geometric interpretation:}

\[
(\nabla \mathbf{v})_f \cdot \mathbf{n}_f
=
\frac{\partial \mathbf{v}}{\partial n}
\Big|_f
\]

On an orthogonal grid, the face normal direction aligns with the
center-to-center direction, so

\[
\frac{\partial \mathbf{v}}{\partial n}
\Big|_f
\approx
\frac{\mathbf{v}_N - \mathbf{v}_P}{d_{PN}}.
\]

\smallskip
\emph{This is a directional derivative in the normal direction.}
\end{frame}




%==================== Slide 13 ====================
\begin{frame}{Assembling the face contributions (one cell)}
For a Cartesian cell $P=(i,j)$, the semi-discrete predictor equation is

\[
\rho A_P \frac{\mathbf{v}_P^{*}-\mathbf{v}_P^{n}}{\Delta t}
+
\sum_{f\in\{E,W,N,S\}}
\Big(
\rho\,\mathbf{v}_f\,\phi_f
-
\mu\,(\nabla\mathbf{v})_f\cdot\mathbf{n}_f\,\ell_f
\Big)
=
\mathbf{0}.
\]

\bigskip
Write the sum explicitly:

\[
(\text{E}) + (\text{W}) + (\text{N}) + (\text{S}).
\]

Each face contributes:

\[
\text{Face } f:
\quad
\rho\,\mathbf{v}_f\,\phi_f
-
\mu
\frac{\mathbf{v}_N-\mathbf{v}_P}{d_{PN}}
\,\ell_f.
\]

\bigskip
The update for $\mathbf{v}_P$ is simply
\[
\text{storage} + \text{E} + \text{W} + \text{N} + \text{S} = 0.
\]
\end{frame}

%==================== Slide 14 ====================
%==================== Slide 14 ====================
\begin{frame}{Diffusion term: East and West faces}

Consider only the diffusive contributions
\[
\sum_{f\in\{E,W,N,S\}}
-
\mu\,(\nabla\mathbf{v})_f\cdot\mathbf{n}_f\,\ell_f.
\]

\bigskip
\textbf{East face}
\[
(\nabla\mathbf{v})_E\cdot\mathbf{n}_E
\approx
\frac{\mathbf{v}_{i+1,j}-\mathbf{v}_{i,j}}{\Delta x},
\qquad
\ell_E=\Delta y
\]

\textbf{West face}
\[
(\nabla\mathbf{v})_W\cdot\mathbf{n}_W
\approx
\frac{\mathbf{v}_{i,j}-\mathbf{v}_{i-1,j}}{\Delta x},
\qquad
\ell_W=\Delta y
\]

\bigskip
So E+W contribution is

\[
-\mu \Delta y
\left(
\frac{\mathbf{v}_{i+1,j}-\mathbf{v}_{i,j}}{\Delta x}
-
\frac{\mathbf{v}_{i,j}-\mathbf{v}_{i-1,j}}{\Delta x}
\right).
\]

\end{frame}
%==================== Slide 15 ====================
\begin{frame}{Diffusion term: Simplifying and adding N/S}

Simplify East + West:

\[
-\mu \Delta y
\frac{\mathbf{v}_{i+1,j}-2\mathbf{v}_{i,j}+\mathbf{v}_{i-1,j}}{\Delta x}.
\]

\bigskip
\textbf{North and South faces}

\[
(\nabla\mathbf{v})_N\cdot\mathbf{n}_N
\approx
\frac{\mathbf{v}_{i,j+1}-\mathbf{v}_{i,j}}{\Delta y},
\qquad
\ell_N=\Delta x
\]

\[
(\nabla\mathbf{v})_S\cdot\mathbf{n}_S
\approx
\frac{\mathbf{v}_{i,j}-\mathbf{v}_{i,j-1}}{\Delta y},
\qquad
\ell_S=\Delta x
\]

\bigskip
So N+S gives

\[
-\mu \Delta x
\frac{\mathbf{v}_{i,j+1}-2\mathbf{v}_{i,j}+\mathbf{v}_{i,j-1}}{\Delta y}.
\]

\end{frame}

%==================== Slide 16 ====================
\begin{frame}{Full 2D diffusion term: Laplacian emerges}

Combine all four faces:

\[
\sum_f \text{diffusion}
=
-\mu \left[
\Delta y
\frac{\mathbf{v}_{i+1,j}-2\mathbf{v}_{i,j}+\mathbf{v}_{i-1,j}}{\Delta x}
+
\Delta x
\frac{\mathbf{v}_{i,j+1}-2\mathbf{v}_{i,j}+\mathbf{v}_{i,j-1}}{\Delta y}
\right].
\]

\bigskip
Divide by cell area $A_P=\Delta x\Delta y$:

\[
\boxed{
\mu
\left(
\frac{\mathbf{v}_{i+1,j}-2\mathbf{v}_{i,j}+\mathbf{v}_{i-1,j}}{\Delta x^2}
+
\frac{\mathbf{v}_{i,j+1}-2\mathbf{v}_{i,j}+\mathbf{v}_{i,j-1}}{\Delta y^2}
\right)
}
\]

\bigskip
\centering
\textbf{The discrete Laplacian appears naturally from face fluxes.}

\end{frame}
%==================== Slide 17 ====================
\begin{frame}{Predictor step: full momentum update (no pressure)}

We now assemble the predictor equation
(with pressure temporarily omitted):

\[
\rho A_P \frac{\mathbf{v}_P^{*}-\mathbf{v}_P^{n}}{\Delta t}
=
-\sum_{f}
\rho\,\mathbf{v}_f\,\phi_f
+
\mu \nabla_h^2 \mathbf{v}_P.
\]

\bigskip
On a Cartesian grid:

\[
\nabla_h^2 \mathbf{v}_P
=
\frac{\mathbf{v}_{i+1,j}-2\mathbf{v}_{i,j}+\mathbf{v}_{i-1,j}}{\Delta x^2}
+
\frac{\mathbf{v}_{i,j+1}-2\mathbf{v}_{i,j}+\mathbf{v}_{i,j-1}}{\Delta y^2}.
\]

\bigskip
This produces a provisional velocity
\[
\mathbf{v}^*
\]
that generally does \emph{not} satisfy
\[
\sum_f \phi_f^* = 0.
\]

\end{frame}

\begin{frame}
All face quantities are evaluated using $\mathbf{v}^n$:
\[
\mathbf{v}_f = \mathbf{v}_f^n,
\qquad
\phi_f = (\mathbf{v}_f^n\cdot\mathbf{n}_f)\,\ell_f.
\]

\end{frame}


%==================== Slide 18 ====================

%==================== Slide 20 ====================
\begin{frame}{Pressure correction: from $\mathbf{v}^*$ to $\mathbf{v}^{n+1}$}
\bigskip
Assume we have computed the provisional velocity $\mathbf{v}^*$ from the predictor step.
In general, $\mathbf{v}^*$ is not divergence-free.


We define the corrected velocity by
\[
\boxed{
\mathbf{v}^{n+1}
=
\mathbf{v}^*
-
\frac{\Delta t}{\rho}\nabla p^{n+1}.
}
\]

\bigskip
Take the divergence of both sides:
\[
\nabla\cdot\mathbf{v}^{n+1}
=
\nabla\cdot\mathbf{v}^*
-
\frac{\Delta t}{\rho}\nabla\cdot(\nabla p^{n+1}).
\]

\bigskip
Impose incompressibility at the new time level:
\[
\nabla\cdot\mathbf{v}^{n+1}=0.
\]
\end{frame}


%==================== Slide 21 ====================
\begin{frame}{Pressure Poisson equation}

From incompressibility we obtain

\[
\boxed{
\nabla^2 p^{n+1}
=
\frac{\rho}{\Delta t}\,\nabla\cdot\mathbf{v}^*.
}
\]

\bigskip
\textbf{Important:}

\begin{itemize}
\item This is a \emph{Poisson equation} for $p^{n+1}$.
\item The right-hand side is completely known.
\item It is computed from the provisional velocity $\mathbf{v}^*$.
\item The equation is linear in $p^{n+1}$.
\end{itemize}

\bigskip
After solving for $p^{n+1}$, we update

\[
\mathbf{v}^{n+1}
=
\mathbf{v}^*
-
\frac{\Delta t}{\rho}\nabla p^{n+1}.
\]

\end{frame}


%==============================================================

%==================== Final Slide 1 ====================
\begin{frame}{Algorithm Summary: Projection Method (FV)}

At each time step:

\bigskip
\textbf{1. Predictor (momentum without pressure)}
\[
\mathbf{v}^* = 
\mathbf{v}^n
+ \Delta t \; \mathcal{N}(\mathbf{v}^n)
+ \Delta t \; \mathcal{D}(\mathbf{v}^n)
\]

\bigskip
\textbf{2. Solve pressure Poisson equation}
\[
\nabla^2 p^{n+1}
=
\frac{\rho}{\Delta t}\nabla\cdot\mathbf{v}^*
\]

\bigskip
\textbf{3. Correct velocity}
\[
\mathbf{v}^{n+1}
=
\mathbf{v}^*
-
\frac{\Delta t}{\rho}\nabla p^{n+1}
\]

\bigskip
\textbf{4. Enforce boundary conditions}

\end{frame}


%==================== Final Slide 2 ====================
\begin{frame}{What did we just do?}

\begin{itemize}
\item Converted Navier–Stokes into conservative flux form.
\item Applied the finite-volume method.
\item Approximated face fluxes on an orthogonal grid.
\item Derived the discrete Laplacian naturally.
\item Enforced incompressibility via a Poisson solve.
\end{itemize}

\bigskip
\centering
\Large
Pressure enforces mass conservation.
\end{frame}


%==================== Final Slide 3 ====================
\begin{frame}{Next: Applications}

\begin{itemize}
\item Lid-driven cavity (Python implementation)
\item Flow around a cylinder (OpenFOAM)
\item Grid effects and stability
\item Linear solvers for Poisson equation
\end{itemize}

\bigskip
\centering
\Large
The mathematics is complete.  
Now we compute.
\end{frame}

%==========================================================










\end{document}